% .claude/skills/md-to-pdf/assets/templates/thesis/main.tex — 学位論文の雛形
% LuaLaTeX + jlreq(book) + biblatex/biber
%
% 使い方:
%   1. このフォルダごと work/<プロジェクト>/drafts/<論文名>/ にコピーする
%   2. latexmk main.tex  （VSCode の LaTeX Workshop なら保存時に自動ビルド）
%   パスの書き換えは不要。ただし雛形の置き場所のままビルドすると、スキルの中に出力が散らかる。
%
% 章は chapters/ に1ファイルずつ置き、ここから \input で読み込む。
% 1ファイルが長くなりすぎないので、章単位で AI に渡しやすい。

\documentclass[a4paper,12pt,oneside,openany,book]{jlreq}

% 共通プリアンブルと共通の書誌データ。パスは書かない（.latexmkrc が解決する）
\input{preamble-ja}
\addbibresource{references.bib}

\title{博士論文のタイトル\\
  {\large サブタイトルをここに置く}}
\author{氏名}
\date{\today}

\begin{document}

\maketitle

% ===== 前付（目次など。ページ番号はローマ数字） =====
\frontmatter
\tableofcontents

% ===== 本文 =====
\mainmatter
% chapters/1_introduction.tex — 第1章 序論
\chapter{序論}

ここに導入の段落を書く。脚注形式の引用は \verb|\autocite| を使う\autocite{yamada2020}。
ページを指定する場合は \verb|\autocite[10--15]{yamada2020}| のように書く\autocite[10--15]{yamada2020}。

\section{先行研究の到達点と課題}

先行研究を整理し、何が明らかにされ、何が残されているかを示す\autocite[42]{smith2019}。
同じ文献を再度引用すると、短縮形（日本語文献なら「著者、前掲書」）になる\autocite[20]{yamada2020}。

\section{本論文の課題と方法}

問いと、それに答えるための史料・方法を述べる。

\section{本論文の構成}

各章の内容を簡潔に予告する。


% 部（\part）で章をまとめる場合。使わないなら \part の行を消すだけでよい。
\part{第1部の題名}
% chapters/2_chapter.tex — 第2章
\chapter{章の題名}

本文をここに書く。

\section{節の題名}

\subsection{項の題名}

未公刊史料を引用する例\autocite{archive_sample}。刊行史料を引用する例\autocite{printed_sample}。
これらは \verb|references.bib| の \verb|keywords| で史料に分類され、
参考文献一覧では研究文献と別の節に出力される。


\part{第2部の題名}
% chapters/3_chapter.tex — 第3章
\chapter{章の題名}

本文をここに書く。

\section{節の題名}

章ファイルは1章につき1ファイルにしておくと、章単位で AI に読ませたり
書き直させたりしやすい。章を増やすときは、このファイルを複製して
\verb|main.tex| に \verb|\input{chapters/...}| の行を足す。


% 結論を部に属させない場合は、目次上でも前の部から視覚的に切り離す
\addtocontents{toc}{\protect\addvspace{1.5em}}
% chapters/4_conclusion.tex — 結論
\chapter{結論}

各章で論じたことをまとめ、問いに答える。

\section{本論文の知見}

\section{残された課題}


% ===== 後付（参考文献一覧） =====
\backmatter

% \clearpage を無効化して、分割した文献リストを続けて組む
\begingroup
\let\clearpage\relax

\chapter*{参考文献一覧}
\addcontentsline{toc}{chapter}{参考文献一覧}

% 文献の分類は .bib の keywords フィールドで行う。
% 例: keywords = {source-archive} と書いた項目だけが下の「未公刊史料」に出る。
% 分類が不要なら、この節をまるごと消して
%   \printbibliography[title={参考文献}]
% の1行に置き換えればよい。

\section*{I\quad 史料}

\subsection*{1.\quad 未公刊史料}
\printbibliography[heading=subsubbibliography, keyword=source-archive, title={}]

\subsection*{2.\quad 刊行史料}
\printbibliography[heading=subsubbibliography, keyword=source-printed, title={}]

\section*{II\quad 研究文献}

% notkeyword で「史料に分類されなかったもの」を拾う。
% 史料のキーワードを増やしたら、ここにも notkeyword を足すこと。
\subsection*{1.\quad 欧語文献}
\printbibliography[heading=subsubbibliography,
  notkeyword=source-archive, notkeyword=source-printed, notkeyword=japanese,
  title={}]

\subsection*{2.\quad 邦語文献}
\printbibliography[heading=subsubbibliography,
  notkeyword=source-archive, notkeyword=source-printed, keyword=japanese,
  title={}]

\endgroup

\end{document}
